%-------------------------
% Resume in LaTeX
% Author: Suraj Kushwaha
%-------------------------

\documentclass[letterpaper,10pt]{article}

\usepackage{latexsym}
\usepackage[empty]{fullpage}
\usepackage{titlesec}
\usepackage{marvosym}
\usepackage[usenames,dvipsnames]{color}
\usepackage{verbatim}
\usepackage{enumitem}
\usepackage[hidelinks]{hyperref}
\usepackage{fancyhdr}
\usepackage[english]{babel}
\usepackage{tabularx}
\usepackage{fontawesome5}

\pagestyle{fancy}
\fancyhf{}
\fancyfoot{}
\renewcommand{\headrulewidth}{0pt}
\renewcommand{\footrulewidth}{0pt}

% Adjust margins
\addtolength{\oddsidemargin}{-0.5in}
\addtolength{\evensidemargin}{-0.5in}
\addtolength{\textwidth}{1in}
\addtolength{\topmargin}{-.5in}
\addtolength{\textheight}{1.0in}

\urlstyle{same}
\raggedbottom
\raggedright
\setlength{\tabcolsep}{0in}

% Sections formatting
\titleformat{\section}{
  \vspace{-4pt}\scshape\raggedright\large
}{}{0em}{}[\color{black}\titlerule \vspace{-5pt}]

%-------------------------
% Custom commands

\newcommand{\resumeItem}[1]{
  \item\small{{#1 \vspace{-2pt}}}
}

\newcommand{\resumeSubheading}[4]{
  \vspace{-2pt}\item
    \begin{tabular*}{0.97\textwidth}[t]{l@{\extracolsep{\fill}}r}
      \textbf{#1} & #2 \\
      \textit{\small#3} & \textit{\small #4} \\
    \end{tabular*}\vspace{-7pt}
}

\newcommand{\resumeProjectHeading}[2]{
    \item
    \begin{tabular*}{0.97\textwidth}{l@{\extracolsep{\fill}}r}
      \small#1 & #2 \\
    \end{tabular*}\vspace{-7pt}
}

\renewcommand\labelitemii{$\vcenter{\hbox{\tiny$\bullet$}}$}

\newcommand{\resumeSubHeadingListStart}{\begin{itemize}[leftmargin=0.15in, label={}]}
\newcommand{\resumeSubHeadingListEnd}{\end{itemize}}
\newcommand{\resumeItemListStart}{\begin{itemize}}
\newcommand{\resumeItemListEnd}{\end{itemize}\vspace{-5pt}}

%-------------------------------------------

\begin{document}

%----------HEADING----------
\begin{center}
    \textbf{\Huge \scshape Suraj Kushwaha} \\ \vspace{4pt}
    \small
    \faPhone\ +91\,91067\,64917 \enspace $|$ \enspace
    \href{mailto:work@surajkuushwaha.com}{\faEnvelope\ work@surajkuushwaha.com} \enspace $|$ \enspace
    \href{https://github.com/surajkuushwaha}{\faGithub\ surajkuushwaha} \enspace $|$ \enspace
    \href{https://linkedin.com/in/surajkuushwaha}{\faLinkedin\ surajkuushwaha}
\end{center}

%-----------SNAPSHOT-----------
% \section{At a Glance}
% \begin{tabular*}{\textwidth}{l@{\hspace{6pt}}l}
%   \small\textbf{What I build:}   & \small Distributed \textbf{Python}/\textbf{Go} backends, event-driven microservices, and public APIs at production scale \\[4pt]
%   \small\textbf{What I've done:} & \small \textbf{100M+ req/month} at $<$100ms p99 \enspace $\cdot$ \enspace \textbf{35\%} latency cut \enspace $\cdot$ \enspace \textbf{99.98\%} uptime across 30+ deployments \enspace $\cdot$ \enspace \textbf{3$\times$} API adoption \enspace $\cdot$ \enspace \textbf{40\%} MTTR reduction \\
% \end{tabular*}

%-----------SUMMARY-----------
\section{Summary}
\small{
Infrastructure \& platform engineer focused on event-driven architecture and reliability at scale. I build the \textbf{Python}/\textbf{Go} systems behind \textbf{100M+} monthly requests, own public APIs end-to-end, and increasingly ship LLM-backed features (LangChain, embeddings, eval gating) into production.
}

\vspace{4pt}

\vspace{4pt}

\vspace{4pt}

%-----------EXPERIENCE-----------
\section{Work Experience}
  \resumeSubHeadingListStart

    \resumeSubheading
      {Backend Engineer - Infrastructure \& Platform}{July 2022 - Present}
      {CultureX Entertainment Private Limited}{Gujarat, India}
      \resumeItemListStart

        \resumeItem{
          \textbf{Distributed Systems \& Scale:} Reduced p99 API latency by \textbf{35\%} and sustained \textbf{100M+ monthly requests} at sub-100ms response time by redesigning a monolithic pipeline into event-driven microservices using \textbf{Python (Celery)}, \textbf{Go} workers, AWS SQS, EventBridge, and Lambda - with idempotent batch processing and distributed concurrency controls preventing duplicate state under fan-out load.
        }

        \resumeItem{
          \textbf{Public API Product:} Grew external developer adoption \textbf{3$\times$} in two quarters by owning the full lifecycle of a versioned REST API - schema design, rate-limiting, and OpenAPI docs - enabling zero-touch partner integration and reducing inbound engineering support requests to near zero.
        }

        \resumeItem{
          \textbf{Zero-Downtime Migration:} Delivered \textbf{99.98\% uptime} through a 6-month live migration of 30+ global brand tenants from fragmented multi-database silos to a unified cloud-native SaaS architecture on AWS (EC2, RDS, S3), cutting data-consistency incidents by \textbf{80\%} with no customer-facing outage.
        }

        \resumeItem{
          \textbf{Observability:} Reduced MTTR for production incidents by \textbf{40\%} by instrumenting end-to-end distributed tracing, structured JSON logging, and CloudWatch anomaly alerts - shifting the team from reactive firefighting to proactive detection.
        }

        \resumeItem{
          \textbf{AI Feature Delivery:} Shipped content classification and trend-summarization features to \textbf{50+ brand analysts daily} by building \textbf{Python} microservices integrating LangChain, OpenAI, and LangGraph - with LangSmith evals gating each model update before production rollout.
        }

        \resumeItem{
          \textbf{Delivery Velocity:} Launched \textbf{10+ production MVPs} in 18 months by standardizing Docker-based environments, reusable TypeScript service abstractions, and GitHub Actions CI/CD - cutting new-service setup from 2 days to under 2 hours.
        }

        \resumeItem{
          \textbf{Team Leverage:} Sped up code-review cycles by \textbf{30\%} and raised bar on two junior engineers by introducing TypeScript style guides and mandatory performance checkpoints in CI - catching regression before it reached staging.
        }

      \resumeItemListEnd

  \resumeSubHeadingListEnd

%-----------PROJECTS-----------
% \section{Key Projects}
%   \resumeSubHeadingListStart

%     \resumeProjectHeading
%     % don't use terms like Mastermind d
%     % add numbers(like processing 10k creators in x time)
%     % add optimization number(tech side)
%     % how much time AI saved for getting creator from breif 
%       {\textbf{Mastermind} $|$ \emph{AI Creator-Matching Pipeline} $|$ \textit{Python, LangChain, Qdrant, OpenAI, AWS}}{}
%       \resumeItemListStart

%         \resumeItem{
%           \textbf{Pipeline Design:} Achieved a \textbf{20:1 creator funnel ratio} (industry standard) by building an end-to-end multi-modal AI pipeline that ingests a natural-language brand brief, constructs a structured persona schema via LLM, and scores a pool of \textbf{2,000+ creator profiles} down to the top 100-200 matches per query.
%         }

%         \resumeItem{
%           \textbf{Multi-modal Embeddings:} Improved match precision by vectorizing three independent signals per creator - reel audio (transcribed via Whisper), thumbnail visual descriptions, and post captions - using \textbf{OpenAI embedding models} stored in \textbf{Qdrant}, enabling semantic retrieval beyond keyword or follower-count filters.
%         }

%         \resumeItem{
%           \textbf{Async Infrastructure:} Processed creator media at scale by uploading reels to \textbf{AWS S3}, extracting audio asynchronously via \textbf{SQS/SNS}, and hosting ingestion workers on \textbf{AWS Fargate} - decoupling media fetch latency from scoring latency and enabling parallel processing across the full creator pool.
%         }

%         \resumeItem{
%           \textbf{Feedback Loop Orchestration:} Reduced irrelevant recommendations across iterative refinement cycles by designing a stateful user-feedback mechanism using \textbf{AWS Step Functions} - maintaining per-user exclusion context across re-ranking loops so each iteration tightened niche alignment without restarting the full pipeline.
%         }

%       \resumeItemListEnd

%   \resumeSubHeadingListEnd

%-----------PROJECTS-----------
\section{Key Projects}
  \resumeSubHeadingListStart

    \resumeProjectHeading
  {\textbf{AI Creator-Matching Pipeline} $|$ \textit{CultureX} $|$ \textit{TypeScript, LangChain, Qdrant, OpenAI, AWS}}{}
      \resumeItemListStart

        \resumeItem{
          \textbf{End-to-End Impact:} Cut creator shortlisting from \textbf{4 days of manual research to under 30 minutes} by automating the full brief-to-finalist pipeline - a $\sim$190$\times$ speedup that turned a multi-day analyst task into a single query.
        }

        \resumeItem{
          \textbf{Pipeline Design:} Built an end-to-end multi-modal AI pipeline that ingests a natural-language brand brief, constructs a structured persona schema via LLM, and scores a pool of \textbf{10,000+ creator profiles} down to the top 100-200 matches per query. % adjust shortlist size if you want it to reflect a strict 20:1 funnel (~500)
        }

        \resumeItem{
          \textbf{Multi-modal Embeddings:} Improved match precision by vectorizing three independent signals per creator - reel audio (transcribed via Whisper), thumbnail visual descriptions, and post captions - using \textbf{OpenAI embedding models} stored in \textbf{Qdrant}, enabling semantic retrieval beyond keyword or follower-count filters.
        }

        \resumeItem{
          \textbf{Async Infrastructure at Scale:} Processed \textbf{500,000+ media items} (reels, thumbnails, captions) by uploading reels to \textbf{AWS S3}, extracting audio asynchronously via \textbf{SQS/SNS}, and hosting ingestion workers on \textbf{AWS Fargate} - decoupling media-fetch latency from scoring latency to enable parallel processing across the full creator pool.
        }

        \resumeItem{
          \textbf{Feedback Loop Orchestration:} Reduced irrelevant recommendations across iterative refinement cycles by designing a stateful user-feedback mechanism using \textbf{AWS Step Functions} - maintaining per-user exclusion context across re-ranking loops so each iteration tightened niche alignment without restarting the full pipeline.
        }

      \resumeItemListEnd

  \resumeSubHeadingListEnd

%-----------TECHNICAL SKILLS-----------
\section{Technical Skills}
 \begin{itemize}[leftmargin=0.15in, label={}]
    \small{\item{
     \textbf{Primary:} TypeScript/Node.js (Express, Hono), Python (FastAPI), Go  \\
     \textbf{Distributed Systems:} Event-driven architecture, message queues (SQS, SNS, EventBridge), idempotent processing, concurrency control \\
     \textbf{Databases \& Caching:} PostgreSQL, MongoDB, MySQL, Redis \\
     \textbf{Cloud \& Infra:} AWS (Lambda, SQS, SNS, EventBridge, S3, EC2, RDS), Docker, Kubernetes \\
     \textbf{AI/ML Infra:} LangChain, LangGraph, LangSmith, OpenAI API \\
     \textbf{DevOps:} GitHub Actions, CI/CD pipelines, Linux, Git
     \textbf{Observability:} Distributed tracing, structured logging, CloudWatch, anomaly alerting \\
    }}
 \end{itemize}

%-----------EDUCATION-----------
\section{Education}
  \resumeSubHeadingListStart
    \resumeSubheading
      {Bachelor of Engineering - Computer Engineering}{2019 - 2022}
      {Government Engineering College, Modasa}{Gujarat, India}
  \resumeSubHeadingListEnd

\end{document}